%% ==============================================================================
%% main.tex
%% Point d'entrée unique et chef d'orchestre du document
%% Classe : mythesis (bâtie sur KOMA-Script scrbook)
%% ==============================================================================

% Options disponibles pour la classe :
%   - screen : version numérique (marges symétriques, liens cliquables en couleur)
%   - print  : version papier (recto-verso, décalage de pliure BCOR, liens neutres)
\documentclass[screen]{mythesis}

% --- Métadonnées du document ---
\title{Architecture Modulaire pour Documents Scientifiques}
\subtitle{Conception Découplée, Reproductibilité et Typographie Haute Fidélité}
\author{Guillaume \textsc{Auteur}}
\date{17 Octobre 2026}

% Métadonnées institutionnelles (utilisées pour la page de titre et le PDF)
\institution{Institut des Sciences Appliquées \& Université de Recherche}
\faculty{École Doctorale Informatique, Systèmes et Décision}
\department{Laboratoire d'Ingénierie Logicielle et Mathématiques Appliquées}
\degree{Thèse de Doctorat}
\specialty{Informatique, Génie Logiciel \& Typographie Numérique}
\supervisor{Pr. Robert Bringhurst}
\cosupervisor{Dr. Leslie Lamport}
\location{Paris}
\keywords{LaTeX, KOMA-Script, scrbook, architecture logicielle, BibLaTeX, Biber, typographie}

% --- Base bibliographique BibLaTeX ---
\addbibresource{biblio/references.bib}

% ==============================================================================
% Corps du document
% ==============================================================================
\begin{document}

% ------------------------------------------------------------------------------
% 1. Frontmatter : Préliminaires et tables des matières
% ------------------------------------------------------------------------------
\frontmatter

% Page de titre institutionnelle
%% ==============================================================================
%% frontmatter/titlepage.tex
%% Page de titre moderne et épurée utilisant les capacités KOMA-Script
%% ==============================================================================
\begin{titlepage}
  \sffamily
  
  % En-tête : Institution, composante et logo
  \noindent
  \begin{minipage}[t]{0.72\textwidth}
    {\fontsize{13}{16}\bfseries\color{ThesisPrimary}\getinstitution}\par
    \vspace{2mm}
    {\fontsize{10}{13}\color{ThesisSecondary}\getfaculty}\par
    {\fontsize{9}{12}\color{ThesisMuted}\getdepartment}\par
  \end{minipage}%
  \hfill
  \begin{minipage}[t]{0.25\textwidth}
    \raggedleft
    \IfFileExists{figures/logo.pdf}{%
      \includegraphics[height=1.8cm]{logo.pdf}%
    }{%
      \rule{1.8cm}{1.8cm}%
    }
  \end{minipage}

  \vspace{1.5cm}

  % Mention du diplôme et de la discipline
  \begin{center}
    {\fontsize{11}{14}\bfseries\color{ThesisSecondary}\MakeUppercase{\getdegree}}\par
    \vspace{3mm}
    {\fontsize{10}{13}\color{ThesisMuted}Spécialité : \textbf{\getspecialty}}\par
    
    \vspace{1.2cm}

    % Titre principal et sous-titre
    \hrule height 1.5pt \relax
    \vspace{6mm}
    {\fontsize{22}{28}\bfseries\color{ThesisPrimary}\makeatletter\@title\makeatother}\par
    \ifdefined\@subtitle
      \vspace{4mm}
      {\fontsize{14}{18}\color{ThesisSecondary}\makeatletter\@subtitle\makeatother}\par
    \fi
    \vspace{6mm}
    \hrule height 0.5pt \relax

    \vspace{1.2cm}

    % Auteur et date
    {\fontsize{11}{14}\color{ThesisMuted}Présenté et soutenu publiquement par}\par
    \vspace{3mm}
    {\fontsize{15}{18}\bfseries\color{ThesisPrimary}\makeatletter\@author\makeatother}\par
    \vspace{4mm}
    {\fontsize{10}{13}\color{ThesisMuted}Soutenue le \makeatletter\@date\makeatother{} à \getlocation}\par
  \end{center}

  \vfill

  % Composition du jury
  \noindent
  {\small\bfseries\color{ThesisPrimary}Composition du jury :}\par
  \vspace{3mm}
  \noindent
  \begin{tabularx}{\textwidth}{@{} >{\bfseries}p{4.2cm} X >{\raggedleft\arraybackslash}p{4.0cm} @{}}
    \toprule
    Pr. Robert Bringhurst & Professeur, University of British Columbia & Président du jury \\
    Dr. Leslie Lamport    & Directeur de Recherche, INRIA               & Rapporteur \\
    Pr. Donald Knuth      & Professeur Émérite, Stanford University     & Rapporteur \\
    Dr. Claude Shannon    & Chercheur Senior, Bell Labs                 & Examinateur \\
    \getsupervisor        & Professeur, \getinstitution                 & Directeur de thèse \\
    \getcosupervisor      & Maître de Conférences, \getinstitution      & Co-directeur de thèse \\
    \bottomrule
  \end{tabularx}

\end{titlepage}



% Remerciements et résumé bilingue
%% ==============================================================================
%% frontmatter/acknowledgements.tex
%% Remerciements et hommage aux artisans de l'écosystème TeX
%% ==============================================================================
\addchap{Remerciements et Avant-Propos}

Ce manuel a été conçu avec la volonté de démystifier l'architecture des grands projets \LaTeX{}
et d'offrir un modèle de départ pérenne, moderne et immédiatement exploitable.

Mes remerciements les plus sincères vont d'abord aux pionniers de l'informatique typographique :
à Donald Knuth pour avoir imaginé un moteur de composition d'une précision inégalée,
à Leslie Lamport pour avoir rendu ce système accessible aux scientifiques,
ainsi qu'à Markus Kohm, dont la suite KOMA-Script transcende les limites des classes historiques
en apportant la rigueur et l'élégance de la tradition éditoriale européenne.

J'exprime toute ma reconnaissance à l'équipe du projet \LaTeX3,
en particulier à Frank Mittelbach et Ulrike Fischer,
pour leur travail monumental de modernisation du noyau \LaTeX{} contemporain
et l'intégration harmonieuse des nouvelles couches de programmation et de gestion des sockets.

Merci également aux concepteurs de Bib\LaTeX{} et de Biber,
qui ont su libérer la bibliographie académique des entraves de l'ancien format BST,
ouvrant la voie à une personnalisation puissante et transparente.

Puisse ce tutoriel vous faire gagner un temps précieux dans la structuration de vos mémoires,
thèses, rapports d'ingénierie et monographies scientifiques.

%% ==============================================================================
%% frontmatter/abstract.tex
%% Résumés bilingues (Français et Anglais) et mots-clés
%% ==============================================================================
\addchap{Résumé / Abstract}

% --- Résumé en français ---
\section*{Résumé}

La conception de documents scientifiques volumineux, tels que les mémoires de thèse et les rapports d'expertise,
se heurte fréquemment à des difficultés d'évolutivité, de cohérence graphique et de portabilité collaborative.
Ce travail propose une architecture éditoriale modulaire et découplée,
bâtie sur les fondations de la classe \LaTeX{} moderne \code{scrbook} (KOMA-Script).
En instaurant une séparation stricte entre l'orchestration du document maître (\code{main.tex}),
la classe racine (\code{mythesis.cls}),
les profils cibles écran et impression via les fichiers d'options (\code{.clo}),
ainsi que la configuration fine des styles bibliographiques (\code{.bbx}/\code{.cbx}),
cette approche garantit une clarté conceptuelle maximale.
Le système démontre qu'une structure hautement granulaire favorise le travail collaboratif sous contrôle de version,
réduit considérablement la dette technique des préambules monolithiques
et assure une conformité typographique optimale quel que soit le canal de diffusion.

\vspace{0.8em}
\noindent\textbf{Mots-clés :} \getkeywords, classe scrbook, modularité, typographie, Bib\LaTeX.

\vspace{1.5cm}

% --- Abstract en anglais ---
\begin{otherlanguage}{english}
\section*{Abstract}

The authoring of large-scale scientific manuscripts, such as doctoral dissertations and technical monographs,
often suffers from maintainability bottlenecks, visual inconsistencies, and cumbersome collaborative workflows.
This research introduces a decoupled and highly modular \LaTeX{} document architecture,
engineered on top of the modern \code{scrbook} class from the KOMA-Script suite.
By establishing a strict separation of concerns between master document orchestration (\code{main.tex}),
core document class semantics (\code{mythesis.cls}),
target-specific layout options via dedicated class options (\code{.clo}),
and fine-grained citation and bibliography tailoring (\code{.bbx}/\code{.cbx}),
this framework achieves superior transparency and maintainability.
The proposed layout demonstrates that modular decomposition mitigates merge conflicts in version-controlled repositories,
eliminates monolithic preamble pollution,
and delivers publication-grade typographical quality for both digital interactive reading and physical print production.

\vspace{0.8em}
\noindent\textbf{Keywords :} Scientific publishing, KOMA-Script, modular document architecture, \LaTeX{}, Bib\LaTeX.
\end{otherlanguage}



% Sommaire général et listes des illustrations
\tableofcontents
\listoffigures
\listoftables

% ------------------------------------------------------------------------------
% 2. Mainmatter : Corps du document découpé en chapitres
%    L'instruction \include{} est réservée aux chapitres entiers.
% ------------------------------------------------------------------------------
\mainmatter

%% ==============================================================================
%% content/01-introduction.tex
%% Chapitre 1 : Introduction générale et motivation
%% ==============================================================================
\chapter{Introduction Générale}
\label{chap:introduction}

\section{Contexte et Motivations}
\label{sec:intro-contexte}

La rédaction d'un mémoire universitaire ou d'un ouvrage technique de référence constitue une étape charnière dans l'activité scientifique.
Depuis plusieurs décennies, le système typographique \TeX{} initié par \textcite{knuth1984texbook},
complété par l'environnement logiciel \LaTeX{} conçu par \textcite{lamport1994latex},
s'impose comme la référence incontournable de la publication scientifique rigoureuse.
Cependant, à mesure que l'ampleur des documents augmente et que les exigences de diffusion se diversifient,
la gestion d'un fichier source unique ou d'un préambule monolithique présente des limites criantes.

Dans la pratique courante, les auteurs accumulent souvent des centaines de lignes de configurations hétérogènes,
mélangeant la déclaration des polices,
le réglage des marges,
les redéfinitions de commandes mathématiques
et l'inclusion de paquets tiers incompatibles.
Cette approche engendre une fragilité technique majeure :
la moindre modification d'un style ou le changement de moteur de compilation risque d'altérer l'ensemble de la structure.
L'objectif de cette recherche est de formaliser une architecture logicielle pour \LaTeX{}
fondée sur le principe éprouvé de la séparation des préoccupations (\textit{Separation of Concerns}).

\section{Objectifs et Verrous Scientifiques}
\label{sec:intro-objectifs}

Pour répondre à ces enjeux, nous identifions trois verrous méthodologiques essentiels :

\begin{enumerate}
  \item \textbf{Le découplage sémantique :}
  séparer rigoureusement la logique métier (le texte et les équations)
  de la charte graphique globale et de l'orchestration des paquets.
  \item \textbf{La flexibilité du support de diffusion :}
  permettre une transition transparente entre une version optimisée pour la consultation numérique sur écran
  (liens hypertexte interactifs, marges symétriques)
  et une version destinée à l'impression papier de haute qualité
  (compensation de la reliure \code{BCOR}, liens désaturés, neutralité chromatique).
  \item \textbf{La modularité et la traçabilité collaborative :}
  assurer un découpage granulaire des chapitres et des styles
  facilitant le suivi de version sous Git et minimisant les conflits de fusion.
\end{enumerate}

\begin{callout}[Principe d'architecture]
Un document \LaTeX{} de qualité professionnelle doit être abordé comme un projet logiciel :
le fichier maître ordonnance,
la classe encapsule les règles de composition,
les fichiers d'options adaptent le livrable,
et les fichiers de contenu portent la substance intellectuelle.
\end{callout}

\section{Organisation du Document}
\label{sec:intro-plan}

Le présent document est articulé en quatre chapitres principaux :

\begin{itemize}
  \item Le \cref{chap:etat-art} dresse un panorama critique des approches de mise en page académique,
  en s'appuyant notamment sur les préceptes de \textcite{bringhurst2004elements}
  et sur l'écosystème KOMA-Script documenté par \textcite{kohm2024koma}.
  \item Le \cref{chap:methodologie} formalise l'architecture modulaire retenue,
  le mécanisme des fichiers \code{.clo} et la structuration du flux de travail.
  \item Le \cref{chap:resultats} présente les résultats d'évaluation,
  illustrant la robustesse du modèle face à des contraintes typographiques variées.
  \item Enfin, les annexes et la bibliographie rassemblent les compléments d'implémentation.
\end{itemize}


%% ==============================================================================
%% content/02-etat-de-l-art.tex
%% Chapitre 2 : La Classe Racine et les Profils d'Options (.cls & .clo)
%% ==============================================================================
\chapter{La Classe Racine et les Profils d'Options (.cls \& .clo)}
\label{chap:classe-et-options}

\section{Pourquoi Bâtir sur KOMA-Script \code{scrbook} ?}
\label{sec:pourquoi-scrbook}

La distribution standard de \LaTeX{} propose historiquement la classe \code{book}.
Conçue il y a plus de quarante ans, cette classe repose sur des normes typographiques américaines rigides,
avec des marges étroites et un système d'en-têtes peu personnalisable.
Pour modifier son apparence, les utilisateurs sont généralement contraints d'empiler des extensions disparates :
\code{geometry} pour les marges,
\code{fancyhdr} pour les en-têtes,
et \code{titlesec} pour les titres de section.
Or, ces extensions entrent régulièrement en conflit et surchargent les mécanismes internes du moteur.

À l'inverse, la suite KOMA-Script développée par Markus Kohm \autocite{kohm2024koma}
constitue une refonte moderne et cohérente de la mise en page éditoriale.
Fondée sur les règles de proportions harmonieuses de la typographie classique européenne \autocite{bringhurst2004elements,mittelbach2023latex},
la classe \code{scrbook} intègre en son c\oe{}ur :
\begin{itemize}
  \item un calcul dynamique et optique du bloc de texte (\textit{Satzspiegel}) via le paramètre de division \code{DIV},
  \item la prise en compte mathématique de la compensation de reliure (\code{BCOR}),
  \item un système de couches d'en-têtes puissant et natif via \code{scrlayer-scrpage},
  \item et des commandes directes d'ajustement des polices (\code{\textbackslash setkomafont}).
\end{itemize}

\section{Anatomie de la Classe Personnalisée : \code{mythesis.cls}}
\label{sec:anatomie-cls}

Notre classe racine \code{mythesis.cls} encapsule ces fonctionnalités
pour offrir une interface claire et épurée.
Elle prend en charge :

\begin{enumerate}
  \item \textbf{L'interception des options :}
  l'utilisateur peut spécifier \code{[screen]} ou \code{[print]}.
  Toute autre option non reconnue (par exemple la taille de corps \code{12pt} ou l'option \code{draft})
  est automatiquement transmise à la classe parente \code{scrbook} via la macro standard \code{\textbackslash PassOptionsToClass}.
  \item \textbf{La déclaration des métadonnées du document :}
  la classe fournit des accesseurs sémantiques (\code{\textbackslash institution}, \code{\textbackslash degree}, \code{\textbackslash supervisor}, \code{\textbackslash location}, \code{\textbackslash keywords})
  qui alimentent automatiquement la page de garde et les propriétés PDF du visualiseur.
  \item \textbf{La hiérarchie typographique KOMA-Script :}
  les styles des titres sont définis globalement avec \code{\textbackslash setkomafont\{chapter\}} et \code{\textbackslash setkomafont{disposition}},
  garantissant une cohérence visuelle sans modifier le corps du texte.
  \item \textbf{Le chargement conditionnel de l'option de classe (\code{.clo}) :}
  selon le mode choisi, la classe injecte soit \code{mythesis-screen.clo}, soit \code{mythesis-print.clo}.
\end{enumerate}

\section{Le Rôle Déterminant des Fichiers \code{.clo}}
\label{sec:mecanisme-clo}

En \LaTeX{}, l'extension \code{.clo} signifie \textit{Class Option}.
C'est exactement le mécanisme qu'utilise le système \LaTeX{} pour charger les tailles de police de base
(par exemple \code{size10.clo} ou \code{scrsize11pt.clo}).
Dans notre projet, nous détournons élégamment cette fonctionnalité standard
pour définir les profils d'affichage sans alourdir le fichier principal.

\subsection{Le profil écran : \code{mythesis-screen.clo}}
Ce profil est pensé pour la lecture numérique interactive (sur écran d'ordinateur, tablette ou smartphone) :
\begin{itemize}
  \item \textbf{Marges symétriques :} réglage de \code{twoside=false} et \code{BCOR=0mm}.
  Le texte est parfaitement centré sur la page, ce qui offre un confort maximal lors du défilement vertical.
  \item \textbf{Palette chromatique :} couleurs sémantiques (\code{ThesisPrimary},
  \code{ThesisSecondary} et \code{ThesisAccent})
  avec des nuances bleu marine et sarcelle pour un contraste optimal.
  \item \textbf{Hyperliens interactifs :} configuration de \code{hyperref} avec \code{colorlinks=true},
  permettant de distinguer immédiatement les renvois internes (en bleu pétrole), les citations bibliographiques (en vert forêt)
  et les adresses web externes (en violet prune).
  \item \textbf{Signets PDF ouverts :} les panneaux de navigation du lecteur PDF s'ouvrent automatiquement dès l'ouverture du document.
\end{itemize}

\subsection{Le profil impression : \code{mythesis-print.clo}}
Ce profil est optimisé pour la fabrication physique d'un livre ou d'un mémoire relié :
\begin{itemize}
  \item \textbf{Mise en page recto-verso :} réglage de \code{twoside=true} et \code{open=right}.
  Les chapitres majeurs s'ouvrent obligatoirement sur une « belle page » (page impaire située à droite dans le livre ouvert).
  \item \textbf{Compensation de pliure (\code{BCOR}) :} application d'un décalage physique (\code{BCOR=10mm}).
  L'empagement intérieur est décalé pour tenir compte de l'espace absorbé par l'encollage ou la couture de la reliure,
  évitant que le texte ne disparaisse au fond de la gorge du livre.
  \item \textbf{Neutralisation des liens (\code{hidelinks}) :} sur papier, les liens colorés produiraient des trames grisâtres inélégantes.
  L'option \code{hidelinks=true} conserve les liens cliquables dans le PDF numérique sous-jacent,
  mais affiche les textes en noir franc pour une impression laser ou offset impeccable.
\end{itemize}

\section{Tableau Récapitulatif des Deux Profils}
\label{sec:tableau-profils}

Le \cref{tab:comparaison-profils} résume les divergences de comportement entre les deux profils.

\begin{table}[htbp]
  \centering
  \caption{Spécifications comparatives des profils d'options \code{.clo}}
  \label{tab:comparaison-profils}
  \begin{tabularx}{\textwidth}{@{} l >{\raggedright\arraybackslash}X >{\raggedright\arraybackslash}X @{}}
    \toprule
    \textbf{Caractéristique} & \textbf{Profil Écran (\code{mythesis-screen.clo})} & \textbf{Profil Impression (\code{mythesis-print.clo})} \\
    \midrule
    Mode de page           & Recto seul (\code{twoside=false}) & Recto-verso véritable (\code{twoside=true}) \\
    Départ des chapitres   & Page suivante quelconque & Page impaire obligatoire (\code{open=right}) \\
    Compensation reliure   & \qty{0}{\milli\meter} (\code{BCOR=0mm}) & \qty{10}{\milli\meter} (\code{BCOR=10mm}) \\
    Diviseur de grille     & \code{DIV=12} & \code{DIV=12} recalculé avec BCOR \\
    Couleur des liens      & Palette interactive contrastée & Monochrome noir (\code{hidelinks=true}) \\
    En-têtes et folios     & Symétriques (\code{\textbackslash ihead}, \code{\textbackslash ohead}) & Alternés versos/rectos (\code{\textbackslash lehead}, \code{\textbackslash rohead}) \\
    Usage privilégié       & Lecture PDF, soumission en ligne & Reliure thermique, spirale ou cousue \\
    \bottomrule
  \end{tabularx}
\end{table}

\begin{callout}[Recette : Créer un nouveau profil de classe]
Pour ajouter une nouvelle variante (par exemple un profil brouillon \code{mythesis-draft.clo} ou un profil grand format) :
\begin{enumerate}
  \item Créez un fichier \code{mythesis-nom.clo} à la racine.
  \item Déclarez l'option dans \code{mythesis.cls} via \code{\textbackslash DeclareOption\{nom\}\{...\}}.
  \item Activez-le simplement dans \code{main.tex} avec \code{\textbackslash documentclass[nom]\{mythesis\}}.
\end{enumerate}
\end{callout}

%% ==============================================================================
%% content/03-methodologie.tex
%% Chapitre 3 : Styles, Macros Métier et Intégration Graphique (config/ & figures/)
%% ==============================================================================
\chapter{Styles, Macros Métier et Intégration Graphique}
\label{chap:styles-macros-graphisme}

\section{Organisation Modulaire du Dossier \code{config/}}
\label{sec:config-organisation}

Le répertoire \code{config/} sépare rigoureusement deux catégories d'instructions souvent confondues :
les paquets tiers d'une part, et les macros personnelles d'autre part.

\begin{itemize}
  \item \textbf{\code{config/packages.sty} :}
  déclare l'ensemble des dépendances logicielles externes nécessaires au document.
  Il configure leur comportement global (moteur de polices, langues, tableaux, hyperliens)
  dans le respect strict de l'ordre d'importation requis par \LaTeX{}.
  \item \textbf{\code{config/macros.sty} :}
  définit les outils métier spécifiques à la rédaction de votre ouvrage.
  On y trouve vos conventions mathématiques, vos raccourcis textuels
  et la personnalisation des environnements graphiques (\code{tcolorbox}).
\end{itemize}

Grâce à cette scission, si vous souhaitez réutiliser vos macros mathématiques dans un autre article,
il vous suffit d'importer \code{macros.sty} sans devoir importer toute l'infrastructure de la thèse.

\section{L'Art de Configurer les Paquets : \code{packages.sty}}
\label{sec:packages-detail}

Le fichier \code{packages.sty} met en \oe{}uvre plusieurs principes de composition moderne :

\subsection{Détection universelle du moteur typographique}
Le paquet \code{iftex} teste dynamiquement le compilateur utilisé :
\begin{itemize}
  \item Sous \textbf{Lua\LaTeX{}} ou \textbf{Xe\LaTeX{}}, il charge \code{fontspec} et active le crénage natif OpenType.
  \item Sous \textbf{pdf\LaTeX{}}, il charge de manière transparente le codage \code{[T1]\{fontenc\}}, l'encodage \code{[utf8]\{inputenc\}} et la fonte vectorielle \code{lmodern}.
\end{itemize}
Votre document compile ainsi sans friction sur n'importe quelle machine ou forge d'intégration continue.

\subsection{Microtypographie et lisibilité}
Le paquet \code{microtype} applique deux raffinements majeurs de l'imprimerie traditionnelle :
la saillie des marges (\textit{margin protrusion}, permettant aux tirets et à la ponctuation de déborder légèrement dans la marge blanche pour un alignement optique parfait)
et l'expansion subtile des caractères (\textit{font expansion}), qui élimine les césures excessives et les lignes lâches.

\subsection{Tableaux professionnels : la fin des filets verticaux}
Comme l'expliquent \textcite{bringhurst2004elements} et \textcite{mittelbach2023latex},
les tableaux typographiques de haute qualité n'utilisent **jamais** de filets verticaux.
Ces filets fragmentent le regard et gênent la lecture des données.
Notre configuration s'appuie sur le trio moderne :
\begin{enumerate}
  \item \textbf{Filets hiérarchisés :} le paquet \code{booktabs} (\code{\textbackslash toprule}, \code{\textbackslash midrule}, \code{\textbackslash bottomrule}),
  \item \textbf{Largeurs automatiques :} le paquet \code{tabularx} et sa colonne \code{X},
  \item \textbf{Cellules multilignes :} le paquet \code{makecell} pour les retours à la ligne internes.
\end{enumerate}

\section{Macros Sémantiques et Boîtes Graphiques : \code{macros.sty}}
\label{sec:macros-detail}

Le fichier \code{macros.sty} propose une boîte à outils prête à l'emploi.

\subsection{Raccourcis mathématiques rigoureux}
Pour éviter de ressaisir de longues commandes verbeuses, le fichier fournit :
\begin{itemize}
  \item Les vecteurs et matrices en gras droit : \code{\textbackslash vect\{x\}} ($\vect{x}$) et \code{\textbackslash mat\{M\}} ($\mat{M}$).
  \item Les ensembles de nombres usuels : \code{\textbackslash setR} ($\setR$), \code{\textbackslash setN} ($\setN$), \code{\textbackslash setC} ($\setC$).
  \item Les délimiteurs extensibles gérés par \code{mathtools} :
  \code{\textbackslash norm\{x\}} ($\norm{x}$),
  \code{\textbackslash abs\{x\}} ($\abs{x}$)
  et le produit scalaire \code{\textbackslash scal\{x\}\{y\}} ($\scal{x}{y}$).
  \item Le symbole différentiel droit conforme aux normes ISO : \code{\textbackslash diff} ($\diff x$).
\end{itemize}

\subsection{Environnements encadrés avec \code{tcolorbox}}
Trois types d'encadrés personnalisés sont pré-paramétrés :
\begin{enumerate}
  \item \code{\textbackslash begin\{callout\}[Titre]...} : encadré d'information générale ou de remarque méthodologique.
  \item \code{\textbackslash begin\{hypothesis\}[Titre]...} : boîte d'hypothèse scientifique, numérotée par chapitre.
  \item \code{\textbackslash begin\{definitionbox\}[Titre]...} : boîte de définition formelle avec en-tête plein.
\end{enumerate}

Voici une illustration en situation réelle de la boîte d'hypothèse :

\begin{hypothesis}[Découplage architectural strict]
\label{hyp:decouplage}
Tout composant graphique ou typographique peut être modifié ou remplacé
sans qu'aucune ligne des chapitres de contenu (\code{content/*.tex}) ne doive être éditée.
\end{hypothesis}

Grâce à l'intégration de \code{cleveref}, vous pouvez référencer automatiquement cet encadré dans le texte
via \code{\textbackslash cref\{hyp:decouplage\}}, ce qui produit automatiquement : \cref{hyp:decouplage}.

\section{Gestion des Figures Vectorielles : Standalone vs Externalisation TikZ}
\label{sec:figures-tikz}

\subsection{Le problème du recalcul systématique de TikZ}
Par défaut, à chaque compilation du manuscrit, le moteur \LaTeX{} recalcule et redessine l'ensemble
des figures TikZ et des tracés \code{pgfplots} à partir de zéro.
Sur un manuscrit de thèse comprenant 30 à 50 figures élaborées,
ce traitement peut ajouter plusieurs minutes à chaque cycle de compilation,
même si vous avez simplement corrigé une coquille typographique dans un paragraphe textuel.

\subsection{Pourquoi l'externalisation classique plante avec \code{out\_dir = 'build'} ?}
Pour accélérer ce processus, la bibliothèque native \code{\textbackslash usetikzlibrary\{external\}} propose
d'isoler chaque schéma dans un sous-processus d'arrière-plan via l'option \code{-shell-escape},
générant un PDF d'image réutilisé aux compilations ultérieures.
Toutefois, dès que le projet utilise un dossier de sortie isolé (\code{\$out\_dir = 'build'}),
deux conflits majeurs surgissent :
\begin{enumerate}
  \item \textbf{Conflit de répertoire :} TikZ cherche à écrire ses fichiers temporaires et ses sommes de contrôle
  dans un sous-dossier (ex. \code{tikz-cache/main-fig0.md5}). Si ce sous-dossier n'est pas préalablement créé
  à l'intérieur de \code{build/}, la compilation s'interrompt avec l'erreur :
  \code{! I can't write on file 'tikz-cache/main-fig0.md5'}.
  \item \textbf{Conflit de contexte de processus :} le processus principal travaille dans \code{build/},
  mais le sous-processus lancé par \code{\textbackslash write18} s'exécute à la racine.
  Le PDF généré par le sous-processus se retrouve alors égaré à la racine tandis que le moteur principal l'attend dans \code{build/}.
\end{enumerate}

\subsection{La solution reine en thèse : le paquet \code{standalone} (Option A)}
Pour s'affranchir définitivement de ces écueils et respecter les règles de sécurité
(évitant l'activation d'un \code{-shell-escape} non restreint),
notre architecture adopte le paquet \code{standalone}.

Chaque figure devient un fichier \LaTeX{} autonome dans le dossier \code{figures/tikz/} :
\begin{enumerate}
  \item \textbf{Édition ultra-rapide :} vous pouvez ouvrir \code{figures/tikz/schema-principe.tex}
  et le compiler isolément en moins d'une seconde dans votre éditeur pour ajuster vos n\oe{}uds.
  \item \textbf{Inclusion fluide :} dans le chapitre,
  l'instruction \code{\textbackslash includestandalone} extrait le tracé sans réexécuter de préambule :
\end{enumerate}

\begin{figure}[htbp]
  \centering
  \includestandalone[width=\textwidth]{figures/tikz/schema-principe}
  \caption[Schéma d'orchestration modulaire]{Schéma d'orchestration modulaire inclus via le paquet \code{standalone}}
  \label{fig:architecture-projet}
\end{figure}

Comme démontré sur le \cref{fig:architecture-projet},
le schéma hérite fidèlement de la palette chromatique de la thèse (\code{ThesisPrimary}, \code{ThesisSecondary})
grâce à la directive \code{\textbackslash providecolor} déclarée dans l'en-tête du fichier autonome.

\begin{callout}[Alternative : Précompilation en PDF]
Pour un temps de compilation de thèse strictement nul sur les figures stabilisées,
vous pouvez compiler le fichier \code{schema-principe.tex} une fois pour toutes
en \code{schema-principe.pdf}, puis l'appeler classiquement par \code{\textbackslash includegraphics}.
Vous obtenez ainsi une vélocité maximale sans compromettre la source vectorielle.
\end{callout}

\section{Sanctuarisation du Code Source Externe : le Dossier \code{listings/}}
\label{sec:listings-externe}

Une thèse ou un rapport technique intègre fréquemment des implémentations logicielles,
des scripts d'analyse ou des modules algorithmiques.
Le copier-coller de code directement dans les fichiers \code{.tex} constitue une mauvaise pratique notoire :
il détruit l'indentation, corrompt les caractères spéciaux
et empêche d'exécuter ou de tester le code en dehors du document.

Pour préserver une hygiène irréprochable, les codes sources sont placés
dans le répertoire dédié \code{listings/}.
Ils conservent leur extension native (\code{.py}, \code{.c}, \code{.rs}, \code{.jl}),
ce qui permet de les exécuter, de les profiler et de les tester unitairement avec vos outils habituels.
Dans le document \LaTeX{}, le code est inséré sans aucune duplication grâce au paquet \code{listings} :

\lstinputlisting[
  language=Python,
  caption={Pipeline d'optimisation modulaire importé directement depuis \code{listings/algorithm-pipeline.py}},
  label={lst:pipeline-python}
]{listings/algorithm-pipeline.py}

Comme illustré dans le \cref{lst:pipeline-python},
la mise en page hérite automatiquement de la typographie,
de la numérotation des lignes et de la palette chromatique (\code{ThesisPrimary}, \code{ThesisAccent}) définies dans \code{config/macros.sty}.

\section{Insertion Propre de Documents et Schémas PDF Externes}
\label{sec:insertion-pdf}

L'intégration de fichiers au format PDF dans une thèse répond à deux besoins fondamentalement distincts :
l'inclusion d'illustrations vectorielles d'une part,
et l'annexion de documents administratifs ou d'articles scientifiques complets d'autre part.

\subsection{Figures vectorielles : résolution infinie et rognage optique}
Pour les graphiques et schémas issus d'outils externes (Inkscape, Adobe Illustrator, Python ou MATLAB),
le format vectoriel PDF est la règle absolue :
contrairement aux images matricielles (\code{.png}, \code{.jpg}),
il ne pixellise jamais à l'agrandissement et garantit une impression d'une netteté parfaite.

Deux bonnes pratiques indispensables assurent une intégration professionnelle :
\begin{enumerate}
  \item \textbf{Omission de l'extension :} ne précisez jamais l'extension \code{.pdf} dans l'instruction :
  \begin{center}
    \code{\textbackslash includegraphics[width=0.8\textbackslash textwidth]\{figures/mon-schema\}}
  \end{center}
  \LaTeX{} choisira automatiquement le format vectoriel optimal.
  \item \textbf{Élimination des marges blanches parasites :} les logiciels de dessin exportent souvent une page A4 entière autour d'un petit schéma. Deux méthodes permettent d'isoler la boîte englobante exacte :
  \begin{itemize}
    \item \textbf{En amont (recommandé) :} exécuter l'outil en ligne de commande \code{pdfcrop} :
    \begin{center}
      \code{pdfcrop --margins 2 figures/mon-schema.pdf figures/mon-schema.pdf}
    \end{center}
    \item \textbf{Dans \LaTeX{} :} utiliser les clés de rognage \code{trim} et \code{clip} :
    \begin{center}
      \code{\textbackslash includegraphics[trim=10 20 10 20, clip, width=0.7\textbackslash textwidth]\{figures/mon-schema\}}
    \end{center}
  \end{itemize}
\end{enumerate}

\subsection{Documents multipages et publications en annexe : le paquet \code{pdfpages}}
Pour intégrer des pièces jointes officielles (conventions de recherche, attestations de jury, tirés-à-part d'articles publiés),
l'emploi de \code{\textbackslash includegraphics} est inadapté car il ne capture que la première page et tronque le document.

Le paquet \code{pdfpages} (\cref{sec:packages-detail}) répond à ce besoin via la commande \code{\textbackslash includepdf}.
Cependant, son usage direct présente un piège majeur :
par défaut, le style de page appliqué est \code{empty},
ce qui supprime les en-têtes et le numéro de page,
créant une rupture inacceptable dans la pagination continue.

Notre configuration dans \code{macros.sty} propose la commande dédiée :
\begin{center}
  \code{\textbackslash includepdfpaper[pages=\{1-3\}]\{figures/sample-annex-doc.pdf\}}
\end{center}

Cette commande applique deux réglages indispensables :
\begin{itemize}
  \item \code{pagecommand=\{\textbackslash thispagestyle\{scrheadings\}\}} :
  maintient le bandeau supérieur de chapitre et le folio (numéro de page) KOMA-Script au-dessus du document externe.
  \item \code{scale=0.82} :
  applique une réduction d'échelle subtile afin que le document annexé ne chevauche pas les en-têtes et respecte la grille typographique de votre ouvrage.
\end{itemize}

%% ==============================================================================
%% content/04-resultats.tex
%% Chapitre 4 : Résultats expérimentaux et validation
%% ==============================================================================
\chapter{Résultats Expérimentaux et Validation}
\label{chap:resultats}

\section{Protocole Expérimental et Métriques}
\label{sec:resultats-protocole}

Pour valider l'efficience et la robustesse de notre architecture,
nous avons conduit une campagne de bancs d'essai sur plusieurs jeux de données représentatifs.
Les grandeurs physiques et les métriques de performance sont systématiquement formulées
à l'aide du paquet \code{siunitx}, assurant le respect des normes métrologiques internationales.
À titre d'exemple, la puissance crête mesurée est de \qty{142.8}{\kilo\watt},
tandis que la vitesse de convergence moyenne s'établit à \qty{1.45e-3}{\second\per\iter}.

\begin{definitionbox}[Indice d'efficience typographique ($\eta_T$)]
\label{def:efficience}
Soit $N_p$ le nombre total de pages d'un document et $T_c$ le temps de compilation complet avec Biber.
L'indice d'efficience typographique est défini par le ratio :
\begin{equation}
  \eta_T = \frac{N_p}{T_c} \quad \left[\text{pages}\cdot\unit{\second^{-1}}\right]
\end{equation}
Un ratio élevé témoigne d'une compilation optimisée sans goulets d'étranglement dans le préambule.
\end{definitionbox}

\section{Analyse Comparative des Performances}
\label{sec:resultats-analyse}

Les données quantitatives obtenues lors des campagnes de test
sont récapitulées dans le \cref{tab:resultats-synthese}.
Comme préconisé par les bonnes pratiques typographiques,
ce tableau utilise exclusivement les filets horizontaux du paquet \code{booktabs}
(\code{\textbackslash toprule}, \code{\textbackslash midrule}, \code{\textbackslash bottomrule})
et proscrit tout filet vertical afin de favoriser le balayage visuel continu de la donnée.

\begin{table}[htbp]
  \centering
  \caption{Synthèse comparative des performances selon la configuration de compilation}
  \label{tab:resultats-synthese}
  \begin{tabularx}{\textwidth}{@{} l r r r X @{}}
    \toprule
    \textbf{Configuration} & \textbf{Mémoire (\unit{\mega\byte})} & \textbf{Temps (\unit{\second})} & \textbf{Score $\eta_T$} & \textbf{Observations} \\
    \midrule
    Modèle Monolithique   & \num{248.5} & \num{8.42} & \num{12.4} & Avertissements de fontes et collision de paquets \\
    Extensions disparates & \num{310.2} & \num{9.15} & \num{11.5} & Redondance des inclusions de macros \\
    \textbf{Notre classe (Écran)}   & \num{182.4} & \num{4.18} & \num{25.1} & Navigation fluide, hyperliens dynamiques \\
    \textbf{Notre classe (Papier)}  & \num{185.0} & \num{4.22} & \num{24.9} & Géométrie BCOR compensée, contraste maximal \\
    \bottomrule
  \end{tabularx}
\end{table}

Ces résultats corroborent les observations de \textcite{dupont2024methodologie} :
la rationalisation des couches d'importation permet de diviser par deux le temps de compilation
tout en garantissant une totale immunité face aux régressions visuelles.
Comme le démontrent également les travaux de \textcite{goodfellow2016deep} et \textcite{vaswani2017attention},
la rigueur dans la structuration des composants logiciels et des pipelines de traitement
constitue un facteur déterminant pour la reproductibilité expérimentale.

\section{Discussion et Synthèse}
\label{sec:resultats-discussion}

L'évaluation qualitative met en lumière trois bénéfices majeurs pour l'utilisateur :
\begin{itemize}
  \item \textbf{Lisibilité du code source :} un nouvel intervenant identifie instantanément
  où ajouter un chapitre, une macro ou une entrée bibliographique sans risque d'effets de bord.
  \item \textbf{Polyvalence du livrable :} générer la version finale prête pour l'imprimeur
  ou la version numérique pour le jury ne nécessite qu'un seul basculement d'option.
  \item \textbf{Stabilité pérenne :} l'absence d'extensions obsolètes
  garantit une compatibilité totale avec les distributions \TeX{} Live contemporaines et futures.
\end{itemize}



% ------------------------------------------------------------------------------
% 3. Annexes
% ------------------------------------------------------------------------------
\appendix

%% ==============================================================================
%% backmatter/appendix-a.tex
%% Annexe A : Aide-Mémoire & Guide Pratique des Commandes (Cheat Sheet)
%% ==============================================================================
\chapter{Aide-Mémoire et Guide des Commandes}
\label{anx:aide-memoire}

\section{Récapitulatif des Commandes Métier et Macros}
\label{sec:recap-macros}

Les \cref{tab:recap-math,tab:recap-env-outils} rassemblent les raccourcis sémantiques,
les boîtes visuelles et les commandes d'orchestration définis dans \code{config/macros.sty}.

\begin{table}[htbp]
  \centering
  \caption{Aide-mémoire des notations mathématiques et balisage sémantique textuel}
  \label{tab:recap-math}
  \begin{tabularx}{\textwidth}{@{} l l >{\raggedright\arraybackslash}X @{}}
    \toprule
    \textbf{Commande} & \textbf{Rendu visuel} & \textbf{Usage et signification} \\
    \midrule
    \code{\textbackslash vect\{x\}} & $\vect{x}$ & Vecteur en gras droit/italique \\
    \code{\textbackslash mat\{M\}} & $\mat{M}$ & Matrice formelle en gras droit \\
    \code{\textbackslash setR}, \code{\textbackslash setC} & $\setR$, $\setC$ & Ensembles des réels et complexes \\
    \code{\textbackslash norm\{x\}} & $\norm{x}$ & Norme mathématique délimitée \\
    \code{\textbackslash scal\{u\}\{v\}} & $\scal{u}{v}$ & Produit scalaire délimité \\
    \code{\textbackslash diff} & $\diff x$ & Symbole différentiel droit pour intégrales \\
    \code{\textbackslash term\{mot\}} & \term{mot} & Terme technique introduit pour la première fois \\
    \code{\textbackslash code\{snippet\}} & \code{snippet} & Fragment de code ou nom de fichier en ligne \\
    \bottomrule
  \end{tabularx}
\end{table}

\begin{table}[htbp]
  \centering
  \caption{Aide-mémoire des boîtes graphiques, interactivité et orchestration}
  \label{tab:recap-env-outils}
  \begin{tabularx}{\textwidth}{@{} l l >{\raggedright\arraybackslash}X @{}}
    \toprule
    \textbf{Environnement / Commande} & \textbf{Type} & \textbf{Usage et comportement} \\
    \midrule
    \code{callout} & Boîte cadrée & Remarque, conseil ou note méthodologique \\
    \code{hypothesis} & Boîte numérotée & Hypothèse scientifique (compatible \code{cleveref}) \\
    \code{definitionbox} & Boîte d'en-tête & Définition formelle avec bandeau coloré \\
    \code{boiteSynthese} & Boîte synthèse & Synthèse didactique avec cartouche coloré \\
    \code{cadreFigure}   & Cadre arrondi & Encadrement soigné pour figures et schémas \\
    \code{\textbackslash ocgbutton\{id\}\{txt\}} & Bouton cliquable & Bascule interactive d'un calque vectoriel OCG \\
    \code{\textbackslash ocgstack\{vue1\}\{vue2\}} & Double vue & Superposition écran ou linéarisation sur papier \\
    \code{\textbackslash includepdfpaper\{doc\}} & Pages PDF & Insertion multipage avec en-têtes préservés \\
    \code{\textbackslash geometrieLarge} & Marges 1,5 cm & Pleine page pour grands tableaux ou couverture \\
    \code{\textbackslash geometrieWord}  & Marges 2,5 cm & Profil de révision avec marge latérale pour notes \\
    \code{\textbackslash geometrieStandard} & DIV=12 & Restauration de la géométrie canonique KOMA \\
    \code{\textbackslash afficherLogo}, \code{\textbackslash masquerLogo} & Logo scrlayer & Contrôle du calque flottant institutionnel \\
    \code{\textbackslash todoH[st][col]\{n\}\{txt\}} & Note marge & Troncature, infobulle et lien vers la table \\
    \code{\textbackslash missingfigure[opts]\{txt\}} & Réservation & Espace réservé sans bloquer la chaîne \\
    \code{\textbackslash thesistodolist} & Sommaire & Tableau de bord dynamique des tâches de révision \\
    \bottomrule
  \end{tabularx}
\end{table}

\clearpage

\section{Guide des Clés Typographiques KOMA-Script}
\label{sec:koma-font-keys}

Pour adapter l'esthétique générale de votre ouvrage sans toucher au code des chapitres,
vous pouvez modifier les polices attribuées aux éléments KOMA-Script dans \code{mythesis.cls}
via la macro \code{\textbackslash setkomafont\{élément\}\{spécification\}} :

\begin{table}[htbp]
  \centering
  \caption{Principaux éléments typographiques configurables dans \code{mythesis.cls}}
  \label{tab:koma-elements}
  \begin{tabularx}{\textwidth}{@{} l l >{\raggedright\arraybackslash}X @{}}
    \toprule
    \textbf{Élément KOMA} & \textbf{Réglage par défaut} & \textbf{Partie du document affectée} \\
    \midrule
    \code{disposition}      & \code{\textbackslash bfseries\textbackslash sffamily\textbackslash color\{ThesisPrimary\}} & Tous les titres (chapitres, sections) \\
    \code{chapter}          & \code{\textbackslash Huge\textbackslash bfseries\textbackslash sffamily}                   & Titre principal de chapitre \\
    \code{section}          & \code{\textbackslash Large\textbackslash bfseries\textbackslash sffamily}                  & Titre de section de niveau 1 \\
    \code{subsection}       & \code{\textbackslash large\textbackslash bfseries\textbackslash sffamily\textbackslash color\{ThesisSecondary\}} & Titre de sous-section \\
    \code{pageheadfoot}     & \code{\textbackslash small\textbackslash sffamily\textbackslash color\{ThesisMuted\}}     & Texte des en-têtes et pieds de page \\
    \code{pagination}       & \code{\textbackslash small\textbackslash sffamily\textbackslash bfseries}                 & Numéro de page (folio) \\
    \code{captionlabel}     & \code{\textbackslash bfseries\textbackslash sffamily\textbackslash color\{ThesisPrimary\}} & Libellé des légendes (« Figure 1 ») \\
    \code{descriptionlabel} & \code{\textbackslash bfseries\textbackslash sffamily\textbackslash color\{ThesisPrimary\}} & Entrées des listes de description \\
    \bottomrule
  \end{tabularx}
\end{table}

\clearpage

\section{Commandes Terminal Utiles pour la Compilation}
\label{sec:commandes-terminal}

Depuis la racine du dossier \code{mon-projet/} :
\begin{description}
  \item[\code{latexmk}] Compile l'intégralité du document (passes multiples de Lua\LaTeX{} et Biber).
  Le PDF final est placé dans \code{build/main.pdf}.
  \item[\code{latexmk -pvc}] Mode interactif en direct (*continuous preview*) :
  recompile automatiquement le document à chaque sauvegarde de fichier.
  \item[\code{latexmk -c}] Supprime les fichiers auxiliaires de compilation temporaires.
  \item[\code{latexmk -C}] Supprime la totalité des fichiers générés, y compris le PDF produit dans \code{build/}.
\end{description}


% ------------------------------------------------------------------------------
% 4. Backmatter : Glossaire et Bibliographie finale
% ------------------------------------------------------------------------------
\backmatter

%% ==============================================================================
%% backmatter/glossary.tex
%% Glossaire des termes et table des acronymes du projet LaTeX
%% ==============================================================================
\addchap{Glossaire et Sigles}
\label{chap:glossaire}

\section*{Table des Sigles et Acronymes}
\addcontentsline{toc}{section}{Table des Sigles et Acronymes}

\noindent
\begin{tabularx}{\textwidth}{@{} >{\bfseries\color{ThesisPrimary}}p{2.6cm} X >{\raggedleft\arraybackslash\small\color{ThesisMuted}}p{3.5cm} @{}}
  \toprule
  Sigle & Définition / Signification complète & Domaine d'usage \\
  \midrule
  BCOR   & Binding Correction (Compensation de marge pour la reliure)    & Typographie KOMA \\
  BBX    & Bibliography Style File (Mise en page Bib\LaTeX{} des références) & Gestion bibliographique \\
  CBX    & Citation Style File (Mise en page des renvois dans le texte)  & Gestion bibliographique \\
  CLO    & Class Option File (Réglages d'options modulaires de classe)     & Moteur \LaTeX{} \\
  CLS    & Document Class File (Fichier racine de classe de document)     & Moteur \LaTeX{} \\
  CTAN   & Comprehensive \TeX{} Archive Network (Dépôt officiel des paquets) & Écosystème \TeX{} \\
  DIV    & Division Factor (Nombre de bandes pour la grille de page)      & Typographie KOMA \\
  DOI    & Digital Object Identifier (Identifiant pérenne d'article)      & Édition scientifique \\
  KOMA   & KOMA-Script (Ensemble de classes avancées par Markus Kohm)     & Typographie \LaTeX{} \\
  PDF    & Portable Document Format (Standard de diffusion vectoriel)     & Standard numérique \\
  STY    & Package Style File (Fichier d'extension ou de macros)           & Moteur \LaTeX{} \\
  TIKZ   & TikZ ist kein Zeichenprogramm (Moteur de dessin vectoriel)     & Graphisme \LaTeX{} \\
  TOC    & Table of Contents (Table des matières générale)                & Structure éditoriale \\
  \bottomrule
\end{tabularx}

\vspace{1.2cm}

\section*{Glossaire des Notions Typographiques}
\addcontentsline{toc}{section}{Glossaire des Notions Typographiques}

\setkomafont{descriptionlabel}{\bfseries\sffamily\color{ThesisPrimary}}
\begin{description}
  \item[Belle page] Dans un livre imprimé recto-verso, page impaire (située à droite lorsque l'ouvrage est ouvert),
  sur laquelle doivent traditionnellement s'ouvrir les têtes de chapitres (\code{open=right}).
  \item[Compensation de pliure (\code{BCOR})] Élargissement stratégique de la marge intérieure du bloc de texte,
  destiné à compenser l'espace absorbé physiquement par la reliure, l'encollage ou la couture de l'ouvrage.
  \item[Crénage] (\textit{kerning}) : ajustement millimétrique de l'espace optique séparant deux glyphes consécutifs
  (par exemple entre « V » et « A ») afin d'éviter les vides inesthétiques.
  \item[Découplage sémantique] Méthode d'ingénierie consistant à séparer les responsabilités
  d'un projet (orchestration, style, options d'affichage, contenu textuel pur).
  \item[Empagement (\textit{Satzspiegel})] Surface et proportions du rectangle de texte imprimé par rapport aux dimensions
  totales de la feuille de papier, calculé de façon optique par KOMA-Script via le diviseur \code{DIV}.
  \item[Microtypographie] Ensemble de techniques fines (protrusion des marges et expansion des polices)
  automatisées par le paquet \code{microtype} pour équilibrer la densité de gris typographique de la page.
  \item[Saut de ligne sémantique] Pratique de rédaction sous contrôle de version consistant
  à scinder le code source d'un paragraphe à la fin de chaque phrase ou proposition logique,
  rendant les historiques Git clairs et exempts de faux conflits.
\end{description}


% Impression de la bibliographie avec inclusion dans la table des matières
\printbibliography[heading=bibintoc, title={Références Bibliographiques}]

\end{document}


